\documentclass[9pt,a4paper]{extarticle}
\usepackage{parskip}

% ======================
% Packages
% ======================
\usepackage[utf8]{inputenc}
\usepackage[T1]{fontenc}
\usepackage[brazil]{babel}
\usepackage{amsmath, amssymb}
\usepackage{amsthm}
\usepackage{geometry}
\usepackage{listings}
\usepackage{xcolor}
\usepackage{hyperref}
\usepackage{booktabs}
\usepackage{array}
\usepackage{tabularx}
\usepackage{enumitem}
\usepackage{microtype}
\usepackage{csquotes}
\usepackage{natbib}

\definecolor{important}{RGB}{255, 100, 100}
\definecolor{notice}{RGB}{0, 102, 204}

\newcommand{\impt}[1]{\textcolor{important}{#1}}
\newcommand{\ntc}[1]{\textcolor{notice}{#1}}

\setlist{nosep}
\geometry{margin=2cm}

% ======================
% Listings setup (Python)
% ======================
\definecolor{codebg}{rgb}{0.95,0.95,0.95}
\lstset{
    language=Python,
    backgroundcolor=\color{codebg},
    basicstyle=\ttfamily\footnotesize,
    keywordstyle=\color{blue},
    commentstyle=\color{gray},
    stringstyle=\color{red!70!black},
    showstringspaces=false,
    frame=single,
    breaklines=true,
    inputencoding=utf8
}

% ======================
% Document
% ======================
\begin{document}

\title{Introdução \footnote{Material baseado em \citep{chalmers1993ciencia}}}
\author{BCC502 -- Metodologia Científica para Ciência da Computação}
\date{\today}
\maketitle

\tableofcontents

% ==============================================================
\section{Introdução}
% ==============================================================
\begin{itemize}
    \item O que é ciência?
    \item O que queremos dizer quando usamos a palavra ``científica'' como adjetivo?
    \begin{itemize}
        \item Parece implicar algum tipo de mérito.
        \item Algum tipo especial de confiabilidade.
        \item Por quê?
        \item Existe um ``método científico'' que atribui à ciência estas características? 
    \end{itemize}
    \item Este é o trabalho principal de um filósofo da ciência.
    \item Vamos estudar os principais resultados da filosofia da ciência.
    \begin{itemize}
        \item Perspectiva histórica.
        \item Começamos com as ideias mais rudimentares e vamos expandindo para as ideias mais complexas e modernas.
    \end{itemize}
\end{itemize}


\section{Ciência: Concepção de \textit{senso comum}}

\begin{itemize}
    \item Conhecimento científico é conhecimento provado. 
    \item As teorias científicas são derivadas de maneira rigorosa dos dados adquiridos por observação e experimentação.
    \item A ciência é baseada no que podemos ver, ouvir, tocar, etc.
    \item Opiniões ou preferências pessoais e suposições especulativas não têm lugar na ciência.
    \item A ciência é objetiva.
    \item O conhecimento científico é confiável porque é conhecimento provado objetivamente.
\end{itemize}

Esta concepção se tornou popular no século XVII como consequência da revolução científica da época (Galileu, Newton). 

\begin{itemize}
    \item ``Se quisermos compreender a natureza devemos, devemos consultar a natureza''.
    \item É um erro utilizar obras dos antigos (Aristóteles, Bíblia, ....) como fonte de conhecimento.
    \item A experiência é a fonte do conhecimento.
    \item A ciência é uma estrutura construída sobre fatos.
\end{itemize}

``Não foram tanto as observações e experimentos de galileu que causaram a ruptura com a tradição, mas sua atitude em relação à eles. Para ele, os dados eram tratados como dados, e não relacionados a alguma ideia preconcebida... Os dados da observação poderiam ou não se adequar a um esquema conhecido do universo, mas a coisa mais importante, na opinião de Galileu, era aceitar os dados e construir a teoria para se adequar a eles.'' \citep{anthony1957science}

\section{Indutivismo}

Ciência como conhecimento derivado dos dados da experiência

\subsection{Indutivista ingênuo}

\paragraph{A ciência começa com a observação.}
    \begin{itemize}
        \item O observador científico deve ter órgãos sensitivos normais e inalterados e registra fielmente o que pode ver, ouvir, etc.
        \item Sem preconceitos.
        \item Afirmações (proposições) a respeito do estado do mundo, ou de alguma parte dele, podem ser justificadas de maneira direta pelo uso dos sentidos do observador não-preconceituoso. 
        \item Proposições de observação formam a base a partir da qual leis e teorias científicas devem ser derivadas.
    \end{itemize}


\paragraph{Proposições de observação}
        \begin{itemize}
            \item Exemplos: (i) Hoje temos uma lua cheia. (ii) O sol nasceu no leste.
            \item A verdade de tais observações deve ser estabelecida com cuidadosa observação.
            \item Qualquer observador pode estabelecer ou conferir sua verdade pelo uso direto de seus sentidos.
            \item Observadores podem ver por si mesmos. 
        \end{itemize}

\paragraph{Tipos de afirmação}
    \begin{itemize}
        \item \texttt{Afirmações singulares}: Referem-se a uma ocorrência específica ou a um estado de coisas num lugar específico, num tempo específico. \ntc{Proposições de afirmação são afirmações singulares.}
        \item \texttt{Afirmações universais}: São informações gerais que afirmam coisas sobre as propriedades ou comportamento de algum aspecto do universo. 
        \begin{itemize}
            \item Se referem a todos os eventos de um tipo específico em todos os lugares em todos os tempos. Ex. O Sol nasce no leste.
            \item \impt{As leis que constituem o conhecimento científico são sempre afirmações universais.}
        \end{itemize}
    \end{itemize}

\texttt{Se a ciência é baseada na experiência, então por que meios é possível extrair das afirmações singulares, que resultam da observação, as afirmações universais, que constituem o conhecimento científico?}

\paragraph{Resposta Indutivista}

É legítimo generalizar afirmações universais a partir de obbservações afirmações singulares desde que as condições à seguir sejam satisfeitas:

\begin{enumerate}
    \item O número de proposições de observação (afirmações singulares) que forma a base da generalização dever ser grande;
    \item As observações devem ser repetidas sob ampla variedade de condições. 
    \item Nenhuma proposição de observação deve conflitar com a afirmação universal derivada.
\end{enumerate}

\paragraph{O Princípio Indutivista}

Se um grande número de $A$s foi observado sob uma ampla variedade de condições, e se todos estes as $A$s observados possuiam sem exceção a propriedade $B$, então todos os $As$ têm a propriedade $B$. 

\subsection{Explicações e previsões}

Não basta para a ciência gerar afirmações universais é fundamental que a partir das afirmações universais se gerem explicações e previsões.

O tipo de raciocínio envolvido neste tipo de derivação é o \ntc{raciocínio dedutivo}. 

\paragraph{Raciocínio Dedutivo}

Um argumento dedutivo é \textit{logicamente válido} se, supondo que as premissas  são verdadeiras, a conclusão é necessariamente verdadeira --- isto é, é impossível que as premissas sejam verdadeiras e a conclusão falsa. Note que validade não 
garante que as premissas sejam de fato verdadeiras; um argumento válido com premissas verdadeiras é chamado de \textit{sólido} (\textit{sound}).

\texttt{Exemplo 1}

Premissas: 
\begin{enumerate}
    \item Água razoavemente pura congela a cerca de $0^o$
    \item O radiador do meu carro contém água razoavelmente pura 
\end{enumerate}

Conclusão:
\begin{itemize}
    \item Se a temperatura cair abaixo de $0^o$, a água do radiador do meu carro vai congelar
\end{itemize}


\texttt{Exemplo 2}

Premissas: 
\begin{enumerate}
    \item Todos os gatos têm 5 patas
    \item Oswaldo é meu gato 
\end{enumerate}

Conclusão:
\begin{itemize}
    \item Oswaldo tem 5 patas
\end{itemize}

\paragraph{Padrão do raciocínio dedutivo}

Premissas: 
\begin{enumerate}
    \item Leis e teorias
    \item Condições iniciais 
\end{enumerate}

Conclusão:
\begin{itemize}
    \item Previsões e explicações
\end{itemize}

\impt{No indutivismo, leis, teorias e condições iniciais são obtidas por observação direta e indução. Então previsões e explicações podem ser deduzidas deles.}

\section{A atração do indutivismo ingênuo}

\begin{enumerate}
    \item A maioria das pessoas vai achar a explicação indutivista da ciência adequada
    \item Ela fornece uma explicação formalizada de algumas das impressões popularmente mantidas a respeito da ciência.
    \begin{itemize}
        \item Poder de explicação e previsão
        \item Objetividade 
        \item Confiabilidade superior
    \end{itemize}
\end{enumerate}

\paragraph{Poder de explicar e prever}

\begin{itemize}
    \item Leis e teorias podem ser obtidos por indução a partir de afirmações singulares (proposições de observação).
    \item Condições iniciais são proposições de afirmação e podem ser obtidas por observação direta.
    \item Explicações e previsões podem ser obtidas por dedução.
\end{itemize}

\paragraph{Objetividade}

\begin{itemize}
    \item Observação e raciocínio indutivos são processos objetivos.
    \item Proposições de observação podem ser averiguadas por qualquer observador pelo uso normal dos sentidos. Assim não podem depender do gosto, da opinião ou das esperanças do observador. 
    \item As induções satisfazem ou não as condições prescritas, não é uma questão subjetiva.
\end{itemize}

\paragraph{Confiabilidade}

\begin{enumerate}
    \item Seguem das definições indutivistas sobre observação e indução. 
    \item As proposições de afirmação formam a base da ciência e são seguras e confiáveis 
    \item Essa confiabilidade é transmitida às leis e teorias derivadas, desde que as condições de indução sejam satisfeitas
\end{enumerate}


\bibliographystyle{apalike}   
\bibliography{references}

\end{document}

